\newpage
\section*{Impressum}
\label{impressum}
\pagestyle{cropmarksstyle}

\RaggedRight
{\small
Die FOSSGIS 2017 wird gemeinsam vom FOSSGIS e.V. und der Universität Passau organisiert.

\vspace{0.5em}
\includegraphics[width=0.43\linewidth]{FOSSGIS}
\hfill
\includegraphics[width=0.45\linewidth]{uni-passau}
\hfill

\vspace{0.5em}
\noindent Verantwortlich für den Inhalt:\\
FOSSGIS e.V.\\
Römerweg 5\\
79199 Kirchzarten

\vspace{0.5em}
\noindent Diese Programmheft wurde unter Verwendung von \LaTeX\ und 
anderer freier Software zusammengestellt.\\
Quellcode: github.com/fossgis/booklet2017\\
\noindent Satz und Layout: Michael Reichert\\
Titelgestaltung: Michael Reichert unter Verwendung von \emph{Der Passauer Dom vom "`Fünferlsteg"' aus gesehen} von High Contrast/Wikimedia Commons\\
Geodaten: \includegraphics[height=7pt]{copyright}~Open\-Street\-Map-Mitwirkende, osm.org/copyright\\
Karte auf der Rückseite: openstreetmap-carto, gerendert durch Frederik Ramm\\
Raumplan: Tobias Knerr\\
Icons in den Tabellen: SJJB Management, CC-0\\
Lektorat: Frederik Ramm\\
Fotos: "`UP Mensa"' von Kurt Seebauer (S.~\pageref{social-event}})

\vspace{1em}
\noindent \begin{minipage}[htbp]{0.2\textwidth}
\noindent\includegraphics[width=\linewidth]{cc-by-sa-pdf}
\end{minipage}
\hfill
\begin{minipage}[hbtp]{0.74\textwidth}\RaggedRight
\small
Alle Inhalte dieses Programmhefts unterliegen, sofern nicht anders angegeben, 
der Lizenz \emph{Creative Commons Namensnennung Weitergabe unter gleichen Bedingungen 3.0}.
Logos von Firmen und Organisationen sind hiervon ausgenommen.
\end{minipage}


